\documentclass[11pt]{article}
\usepackage[notes]{yalatex}
% \addbibresource{references.bib}   % uncomment for citations

\title{STAT410: Intro to Probability Theory}
\author{Yuval Amit}
\date{}

\begin{document}
\maketitle

\tableofcontents
\newpage

\section{Probability: The Mathematical Study of Uncertainty}

\begin{defn}
    A \textbf{probability space} $(\Omega,\mc F,P)$ is composed of
    \begin{enumerate}[(i)]
        \item The \textbf{sample space} $\Omega$. It contains all possible outcomes of the ``experiment''.
        \item The \textbf{event space} $\mc F$ consisting of \textbf{events}. An event is a collection of outcomes.
        \item The \textbf{probability measure} $P:\mc F \to \R$.
    \end{enumerate}
\end{defn}

\begin{ex}
    Toss a fair coin. This has the sample space $\Omega := \{H,T\}$, the event space $\mc F = \mb P(\Omega)$, and the probability measure $P :\mc F \to \R$ given by
    \begin{align*}
        P(\{H\}) = 1/2, \quad P(\{T\}) = 1/2, \quad P(\emptyset)=0, \quad P(\{H,T\})=1
    \end{align*}
\end{ex}

\begin{ex}
    Roll a fair dice. This has the sample space $\Omega = \{1,\dots,6\}$, the event space $\mc F = \mb P(\Omega)$. Note that $\mc F$ is too big to succinctly write what the probability measure $P$ does. More on this later.
\end{ex}

\begin{ex}
    Pick a number uniformly at random from $[0,1]$. Here, $\Omega = [0,1]$. We pretend that $\mc F = \mb P([0,1])$, but in reality $\mc F \subsetneq \mb P(\Omega)$. Again, we can't explicitly write what $P$ does.
\end{ex}

\begin{defn}
    We define the following axioms for the event space $\mc F$:
    \begin{enumerate}
        \item $\emptyset \in \mc F$ and $\Omega \in \mc F$.
        \item If $A \in \mc F$, then $A^c \in \mc F$.
        \item If $\{A_i\}$ is a sequence of events in $\mc F$, then $\bigcup A_i \in \mc F$.
    \end{enumerate}
    This is the definition of a $\bm \sigma$\textbf{-algebra}.
\end{defn}

\begin{rmk}
    If $A_1,\dots,A_n \in \mc F$, then $\bigcup_{i=1}^nA_i \in \mc F.$
\end{rmk}

\begin{proof}
    Let $A_i =\emptyset$ for all $i > n$. Then, we know $\bigcup_{i=1}^nA_i= \bigcup _{i=1}^\infty A_i\in \mc F$
\end{proof}

\begin{prop}
    If $\{A_i\}$ is a sequence of events in $\mc F$, then $\bigcap A_i \in \mc F$.
\end{prop}

\begin{proof}
    $$\bigcap _{i=1}^\infty A_i = \paren{\bigcup _{i=1}^\infty A_i^c}^c \in \mc F$$
    by axiom 2, 3, and De Morgan's law.
\end{proof}

\begin{defn}
    We define the following axioms for the probability measure $P$:
    \begin{enumerate}
        \item $P(\emptyset) =0$ and $P(\Omega) = 1$.
        \item For any $A \in \mc F$, $P(A) \geq 0$.
        \item (Countable additivity) If $\{A_i\}$ is a sequence of pairwise-disjoint events, then
        $$P\paren{\bigcup_{i=1}^\infty A_i}=\sum_{i=1}^\infty P(A_i)$$
    \end{enumerate}
\end{defn}

\begin{ex}
    Let $\Omega = \{1,\dots,6\}$ be the sample space for rolling a fair dice, and $\mc F = \P(\Omega)$. Recall that $P(\Omega)=1$, and moreover
    $$\Omega = \{1\} \sqcup \{2\} \sqcup \cdots \sqcup \{6\}.$$
    Thus,
    $$1 = P(\Omega) = \sum_{i=1}^6P(\{i\}).$$
    By the assumption that the dice is fair, we know $P(\{i\}) = P(\{j\})$, hence $P(\{i\}) = 1/6$ for all $i \in \{1,\dots,6\}$.
\end{ex}

\begin{ex}
    Suppose $\Omega = [0,1]$, and for any subinterval $I \subset [0,1]$, $P(I)=$ length of $I$. What is the probability you see a rational number? An irrational number?

    Note that for any $x \in [0,1]$, $P(\{x\})=0$ since $\{x\}$ can be considered a closed interval with length $0$. Thus, since $\Q$ is countable, it follows that $$P(\Q \cap [0,1])=0, \quad P((\R\setminus \Q) \cap [0,1]) = 1$$
\end{ex}

\begin{prop}
    $$P(B^c) = 1-P(B)$$
\end{prop}

\section{Continuity of Probability Measures}

Let $B_1 \supset B_2 \supset \cdots$ be a decreasing sequence of sets. Then, we write
$$\lim_{n \to \infty}B_n = B=\bigcap_{i=1}^\infty B_i,$$
which is the greatest lower bound, or the infimum.
\begin{thm}
    The probability measure $P$ commutes with limits. That is, given an ascending sequence $C_1 \subset C_2 \subset \cdots$,
    $$P\paren{\bigcup_{n=1}^\infty C_n}=P\paren{\lim_{n \to \infty}C_n}=\lim_{n\to \infty}P(C_n)$$
\end{thm}

\begin{proof}
    Let $D_1=C_1$, $D_2 = C_2-C_1$, and generally $D_{i+1}:= C_{i+1}-C_i$. Then, observe
    $$\bigcup _{i=1}^\infty C_i = \bigcup _{i=1}^\infty D_i,$$
    hence
    $$P\paren{\bigcup C_i}=P\paren{\bigcup D_i}.$$
    Moreover, since the $D_i$'s are pairwise disjoint, one shows that
    $$P(D_{i+1})=P(C_{i+1})-P(C_i),$$
    hence
    $$\sum_{i=1}^nP(D_i) = P(C_n).$$
    Thus,
    \begin{align*}
        P\paren{\bigcup_{i=1}^\infty C_i} &=P\paren{\bigcup_{i=1}^\infty D_i}=\sum_{i=1}^\infty P(D_i) 
        \\ &= \lim_{n \to \infty}\sum_{i=1}^nP(D_i) = \lim_{n \to \infty}P(C_n)
    \end{align*}
    as desired.
\end{proof}

\section{Inclusion-Exclusion Principle}

\begin{thm}[Inclusion-Exclusion Principle]
$$P(A\cup B) = P(A) + P(B) - P(A \cap B)$$    
\end{thm}

\begin{proof}
    Observe
    $$P(A \cup B) = P(A \cup (B-A)) = P(A) + P(B-A).$$
    Moreover,
    $$P(B) = P(A\cap B) + P(B-A) \implies P(B-A) = P(B) - P(A \cap B),$$
    thus
    $$P(A \cup B) =  P (A) + P(B) - P(A\cap B).$$
\end{proof}

\begin{cor}
    $$P(A \cup B \cup C) = P(A) + P(B) + P(C) - P(A \cap B) - P(B\cap C) - P(A \cap C) + P (A\cap B \cap C)$$
\end{cor}

\begin{prop}
    \begin{align*}
        P\paren{\bigcup_{i=1}^n A_i} \leq \sum_{i=1}^n P(A_i)
        \\ P\paren{\bigcup_{i=1}^\infty A_i} \leq \sum_{i=1}^\infty P(A_i)
    \end{align*}
\end{prop}

\begin{prop}
    If $A \subset B$, then $P(A) \leq P(B)$.
\end{prop}

\section{Intro to Random Variables}

\begin{defn}
    A \textbf{random variable} is a (measurable) function $X : \Omega \to \R$.
\end{defn}

\begin{defn}
    The \textbf{cumulative distribution function}, abbreviated \textbf{cdf}, is
    $$F_X(a) := P(X \leq a)=P(X \in [-\infty, a]).$$

Now, suppose $X$ takes countably many values, i.e. $\{ X(\omega) \mid \omega \in \Omega\}$ is countable. A \textbf{probability mass function}, abbreviated \textbf{pmt}, is a function $p_X$ such that there's sequence of real numbers $\{x_n\}$ with $p_X(x_n)=P(X = x_n)>0$ for all $n$, and $p_X(x)=P(X=x)=0$ for all other $x$.
\end{defn}

\begin{ex}
    Any experiment with finitely many outcomes $X$ is called a discrete random variable. Then, $P_X(x)$ is the probability mass function of $X$. The set where $P_X(x)>0$ is called the \textbf{support.}

    For any set $B$, we have
    $$P(x\in B) = \sum_{x \in \w{supp}(X)\cap B}p_X(x).$$
\end{ex}

\begin{defn}
    A random variable $X$ is said to be \textbf{(absolutely) continuous} if there exists a function $f_X(x)$ such that $f_X(x) \geq 0$, and for any interval $I=[a,b]$,
    $$P(x \in I) = \int_a^bf_X(x)dx$$
\end{defn}

\section{Conditional Probability}

\begin{que}
    Does knowing $A$ has taken place affect our chances of seeing $B$?
\end{que}

\begin{defn}
    Let $A,B$ be events. Then, the probability of $B$ given $A$ is denoted $P(B|A)$.
\end{defn}

\begin{thm}[][label=thm:cond_intersect]
    $$P(B|A) = \frac{P(A\cap B)}{P(A)}$$
\end{thm}

\begin{proof}
    We can give a frequentist proof with
    $$P(B|A) = \lim_{n \to \infty}\frac{\frac{\#A,B \text{ both take place}}{n}}{\frac{\#A \text{ takes place}}{n}}=\frac{P(A\cap B)}{P(A)}.$$
\end{proof}

\begin{ex}
    An experiment consists of tossing a fair coin, followed by rolling a fair D6 if you see heads, and rolling a fair D4 if you see tails. What is the probability that you see heads and two?

    If $A=$ coin is heads and $B=$ dice is two, then our probability is
    $$P(A\cap B)=P(A)\cdot P(B|A) = \frac{1}2{\cdot \frac{1}{6}}=\frac{1}{12}.$$
\end{ex}

\begin{thm}
    $P'(X):=P(X|A)$ is a probability measure. However, $P''(X):=P(B|X)$ is not a probability measure.
\end{thm}

\begin{thm}[Law of Total Probability]
    Let $A_1,\dots,A_n$ partition $\Omega$. and $P(A_i)>0$ for all $i$. Then,
    $$P(B) = \sum_{i\in I}P(B|A_i)P(A_i)$$
    This also holds for countable partitions.
\end{thm}

\begin{thm}[Bayes' Theorem]
    Suppose $A_1,\dots,A_n$ partition $\Omega$, and $P(A_i)>0$. Suppose further that $B$ is an event with $P(B)>0$. Then, for any $j\in\{1,\dots,n\}$,
    $$P(A_j|B) = \frac{P(B|A_j)P(A_j)}{\sum_{i=1}^nP(B|A_i)P(A_i)},$$
    as desired.
\end{thm}

\begin{proof}
    By \Cref{thm:cond_intersect},
    $$P(A_j|B) = \frac{P(A_j \cap B)}{P(B)},$$
    and by the law of total probability
    $$\frac{P(A_j \cap B)}{P(B)}=\frac{P(A_j \cap B)}{\sum_{i=1}^nP(B|A_i)P(A_i)}=\frac{P(B|A_j)\cdot P(A_j)}{\sum_{i=1}^nP(B|A_i)P(A_i)}$$
\end{proof}

\begin{ex}
    The prevalence of a particular disease is $1\%$. A rapid test for this disease has a $2\%$ false positive rate, and a $3\%$ false negative rate. If the test comes up positive, how likely is it that you have the disease?
\end{ex}

\begin{solution}
Let $D$ be the event that you have the disease, and $T$ be the event that you tested positive. Then,
\begin{align*}
    P(D) &= 0.01,
    \\ P(T |D^c)&=0.02,
    \\ P(T^c|D) &= 0.03.
\end{align*}
By Bayes' Theorem,
$$P(D|T) = \frac{P(T|D)P(D)}{P(T|D)P(D) + P(T|D^c)P(D^c)}=\frac{(1-0.03)\cdot 0.01}{(1-0.03)\cdot 0.01 + 0.02\cdot (1-0.01)}=0.3288$$
\end{solution}

\section{Independence}

\begin{defn}
    Events $A$ and $B$ are said to be \textbf{independent} of $P(A\cap B)=P(A)P(B)$.
\end{defn}

\begin{thm}
    Let $A,B$ be events with $P(A),P(B)>0$. If they are independent, they cannot be disjoint.
\end{thm}

\begin{proof}
    $P(A\cap B)=P(A)P(B)>0$.
\end{proof}

\begin{thm}
    For any events $A_1,\dots,A_n$ with $P(A_i)>0$,
    $$P\paren{\bigcap_{i=1}^nA_i}=P\paren{A_n |\bigcap_{i=1}^nA_i}P\paren{A_{n-1}|\bigcap_{i=1}^{n-2}A_i}\cdots P(A_1).$$
\end{thm}

\begin{ex}
    An experiment consists of tossing a fair coin until you see heads for the first time, at which point the experiment terminates. Let $A_i$ be the event you see tails on the $i$'th toss. Are the $A_i$'s independent? What is the probability we see heads for the first time on the $n$'th toss?

    No, the $A_i$'s are not independent. Next, let $B_n$ be the event that heads is seen on the $n$'th toss. Then,
    $$B_n = A_1 \cap A_2 \cap \cdots \cap A_{n-1}\cap A_n^c.$$
    Therefore,
    $$P(B_n) = P\paren{A_n^c |\bigcap_{i=1}^{n-1}A_i} \cdot \prod_{k=1}^{n-1}P\paren{A_k|\bigcap_{i=1}^{k-1}A_i}=(1/2)^n.$$
\end{ex}

\begin{exercise}
    If $A,B$ are independent, show that $A,B^c$ are independent.
\end{exercise}

\section{Random Variables}

\begin{defn}
    For a random variable $X$, and a (measurable) $A \subset \R$, we define the distribution $D_X$ to be
    $$D_X(A) = P(X\in A)$$
\end{defn}

Now, given a random variable $X:\Omega \to \R$, we have
\[ D_X(A)= P (X \in A) = P\paren{\brax{\omega \setbar X(\omega) \in \R}},\]
and
\[F_X(a) = P(X \leq a) = P(X \in (-\infty, a]).\]
Clearly, $D_X$ determins $F_X$, but in fact we'll see that $F_X$ determines $D_X$.

\begin{ex}
    Let $X$ be the outcome of rolling a fair die. Then,
    \[ F_X(a) = \case{0 & a<1 \\ 1/6 & 1 \leq a < 2 \\ 2/6 & 2 \leq a < 3 \\ 
    \vdots \\ 5/6 & 5 \leq a < 6 \\ 1& 6 \leq a. } \]
    In other words,
    \[ F_X(a) = \case{0 & a<1 \\ 1 & a \geq 6 \\ \frac{\floor{a}}{6} & 1 \leq a < 6.} \]
\end{ex}

\begin{prop}
The function $F_X:\R \to \R$ satisfies the following properties:
\begin{enumerate}
    \item $\lim_{X \to \infty} F_X(a) = 1$ and $\lim_{X \to -\infty}F_X(a)=0$
    \item $F_X$ is monotonically increasing
    \item $F_X$ is right continuous, i.e. \[ \lim_{x \to a^+}F_X(x)=F_X(a) \]
\end{enumerate}
\end{prop}

\begin{defn}[Discrete Random Variables]
    $X$ is called \textbf{discrete} if there is a countable $A \subset \R$ such that $P(X \in A)=1$.

    Then, the function $a \mapsto P(X=a)$ is called the \textbf{probability mass function} (pmf), and the probabilites of interest are computed by taking a suitable sum of the pmf.
\end{defn}

\begin{defn}[Continuous Random Variables]
    $X$ is called \textbf{(absolutely) continuous} if there is some function $f_X(x)$ such that $f_X(x) \geq 0$, and for any interval $I$ with endpoints $-\infty \leq a \leq b \leq \infty$,
    \[
    P(X \in I) = \int_a^b f_X(x)dx
    \]
    Note that we then have
    \[
    F_X(a) = \int_{-\infty}^a f_X(x)dx
    \]
\end{defn}

\begin{defn}[Identical Distributions]
    Random variables $X,Y$ are said to be \textbf{identically distributed} if $D_X = D_Y$, which holds if and only if $F_X = F_Y$.
\end{defn}

A good way to think about identical distributions is with an example: how do you generate a random integer in $\{0,1\}$?

One way to do it is flipping a coin, taking $0$ if the coin is tails, and $1$ if the coin is heads. Another way to do it is generating a random real number in $[0,1]$, taking $0$ if it's $<1/2$, and $1$ if it's $\geq1/2$. 

Let $X$ be the coin method and $Y$ be the real number method. Note that $X$ and $Y$ have completely different sample spaces:
\begin{table}[h]
    \centering
    \begin{tabular}{ll}
        $X:\{H,T\} \to \R$ & $X(\{H\})=1, X(\{T\})=0$ \\
        $Y:[0,1]\to \R$ & $Y(\omega)=\case{0 & \omega < 0.5 \\ 1 & \omega \geq 0.5.}$
    \end{tabular}
\end{table}

However, observe
\[
F_X(a) =F_Y(a) = \case{ 0 & a < 0 \\ 1/2 & 0 \leq a < 1 \\ 1 & a \geq 1. }
\]
Thus, $X$ and $Y$ are identically distributed, even though they're completely different (they have different sample spaces!).

Moreover, let $W=1-X$, i.e. $W=1$ on tails and $W=0$ on heads. Then, $X$ and $W$ are still identically distributed!

\begin{rmk}
    \begin{enumerate}
        \item If $F_X$ is piecewise constant, then $X$ is a discrete random variable. In this case, $P(X=a)$ can be computed by $F(a) - \lim_{X \to a^{-}} F_X(a)$
        \item If $X$ is (absolutely) continuous, $F_X$ is continuous
        \item If $F_X$ is differentiable with only finitely many exceptions, then $F_X'(x)=f_X(x)$ for all points where $F_X$ is differentiable. At other points, we may define $f_X$ arbitrarily.
    \end{enumerate}
\end{rmk}

\begin{ex}
Suppose $X$ is a continuous random variable with pdf $f_X(x)=1/2$ for $x \in [-1,1]$.
\begin{enumerate}
    \item[(i)] Compute $P(X>0)$ 
\end{enumerate}
We know
\[
P(X>0)=\int_0^\infty f_X(x)dx = \int_0^1 \frac{1}{2}dx=1/2
\]
\begin{enumerate}
    \item[(ii)] Compute $P(X>1/2 | X>0)$.
\end{enumerate}
Using conditional probability, we see that
\[
P(X>1/2 | X>0)=\frac{P(X>1/2 \cap X>0)}{P(X>0)}=\frac{1/4}{1/2}=\frac{1}{2}.
\]
\begin{enumerate}
    \item[(iii)] Compute the pdf of $X^2$.
\end{enumerate}
Observe that
\[
P(X^2\leq a)=P(-\sqrt a \leq X \leq \sqrt a)=\int_{-\sqrt a}^{\sqrt a}f_X(x)dx.
\]
When $0 \leq a < 1$, this evaluates to $\sqrt a$. Thus, we see
\[
F_{X^2}(a) = \case{0 & a<0 \\ \sqrt a & 0 \leq a < 1 \\ 1 & a \geq 1.}
\]
Note that this is the cdf of $X^2$. Differentiating, we get
\[
f_{X^2}(a) = \case{0 & a<0 \\ \frac{1}{2 \sqrt a} & 0 \leq a <1 \\ 0 & a \geq 1.}
\]
\end{ex}

\begin{defn}
    We say that a random variable $X$ has a \textbf{uniform distribution} on $[a,b]$, denoted $X \sim u[a,b]$, if \[
    f_X(x)=\frac{1}{b-a}
    \]
    for $x \in a,b$.
\end{defn}

\begin{ex}
Suppose $X \sim u[0,1]$. Find the pdf of $Y:=e^X$.
\end{ex}

\begin{solution}
    If $a<1$, then $e^X$ doesn't attain $a$ on $[0,1]$, so $F_Y(a)=0$ for $a<1$. Similarly, when $a>e$, we know with certainty that $e^X<a$, so $F_Y(a)=1$. Finally, if $a \in [1,e]$, then
    \[
    F_Y(a)=P(Y \leq a)=P(e^X \leq a) = P(X \leq \log  a)=P(0 \leq X \leq \log a) =\log a.
    \]
    Thus,
    \[
    F_Y(a) = \case{0 & a < 1 \\ \log (a) & 1 \leq a \leq e \\ 1 & a>e,}
    \]
    and differentiating yields
    \[
    f_Y(a)= \case{0 & a \leq 1 \\ \frac{1}{a} & 1 < a < e \\ 0  & a \geq e.}
    \]
\end{solution}

\section{Expected Values}


\begin{defn}[Expected Value (Discrete)]
    Let $X$ be a discrete random variable with support $A$. Then, the \textbf{expected value} of $X$, denoted $E[X]$, is
        \[
    E[X] = \sum_{x \in A} xp_X(x)=\sum_{x \in A} xP(X=x)
    \]
\end{defn}


\begin{defn}[Expected Value (Continuous)]
    Let $X$ be a continuous random variable. Then, the \textbf{expected value} of $X$, denoted $E[X]$, is
    \[
    E[X] = \int_{-\infty}^\infty xf_X(x) dx
    \]
\end{defn}

\begin{defn}
    We say that the expected value $E[X]$ \textbf{exists} if $E[|X|]< \infty$.
\end{defn}

\begin{ex}
Suppose $X \sim u[0,1]$. Compute $E[X]$ and $E[e^X]$.
\end{ex}

\begin{solution}
    By definition,
    \[
    E[X] = \int_{-\infty}^\infty xf_X(x)dx = \int_0^1xdx = \frac{1}{2}.
    \]
    For $E[e^X]$, one method involves using the pdf we computed earlier, which gives
    \[
    E[e^X]=\int_{-\infty}^\infty yf_{Y}(y)dy = \int_1^e \frac{y}{y}dy=e-1.
    \]
    However, we may use the following general result:
    \[
    E[g(x)]=\int_{-\infty}^\infty g(x)f_X(x)dx,
    \]
    which gives
    \[
    E[e^X]=\int_{0}^1 e^{x}dx=e-1.
    \]
\end{solution}

\end{document}
